\section{Microgrid Planning Under Annual Carbon Budgets and Scenario Reduction}

\ifdefined\TenAPaperMode
\else
This section records the technical formulation, algorithmic design, and validation evidence for the corrected split-cost microgrid-planning study. The paper-style title, abstract, and introduction have been separated into the standalone manuscript \texttt{10a\_paper\_main.tex}; the notebook version below now focuses on the model, theory, and experiment details used to support that paper narrative.
\fi

\subsection{Study Object and Assumptions}

Consider a microgrid over a standard year. Let \(\mathcal{D}=\{1,\dots,365\}\) denote the day set, \(\mathcal{T}=\{1,\dots,T\}\) the intra-day time set, and \(\pi_d\) the empirical probability of day \(d\), with \(\sum_{d\in\mathcal{D}} \pi_d = 1\). In the present application, \(T=24\) and \(\Delta t=1\) hour.

The first-stage decision vector is denoted by
\begin{align}
x := \big(P^{\mathrm{pv}}, P^{\mathrm{bess}}, E^{\mathrm{bess}}, P^{\mathrm{grid}}, \zeta^{\mathrm{CO_2}}, \zeta^{\mathrm{rel}}\big),
\end{align}
where the components represent PV capacity, battery power capacity, battery energy capacity, grid-interconnection capacity, carbon-budget slack, and reliability slack, respectively. When carbon and reliability are treated as hard constraints, one sets \(\zeta^{\mathrm{CO_2}}=\zeta^{\mathrm{rel}}=0\).

For each day \(d\), let \(\xi_d\) denote the scenario data, including base-load profile, flexible-load baseline profile, PV availability, import price, export remuneration, and operating-state information. The corresponding second-stage decision vector is denoted by
\begin{align}
y_d := \big(p^{\mathrm{buy}}_{t,d}, p^{\mathrm{sell}}_{t,d}, p^{\mathrm{pv}}_{t,d}, p^{\mathrm{ch}}_{t,d}, p^{\mathrm{dis}}_{t,d}, e_{t,d}, p^{\mathrm{ls}}_{t,d}, \ell^{\mathrm{flex}}_{t,d}, u^+_{t,d}, u^-_{t,d}, I_{t,d}, z^{(1)}_{t,d}, z^{(2)}_{t,d}\big)_{t\in\mathcal{T}}.
\end{align}
When the low-SOC degradation regularizer is disabled, the auxiliary variables \(z^{(1)}_{t,d},z^{(2)}_{t,d}\) and the associated coefficients are simply omitted from the realized recourse model.

The model is built under the following assumptions.

\begin{assumption}[Exogenous day data]
For each day \(d\), the scenario vector \(\xi_d\) is exogenous and contains all uncertain quantities that affect recourse feasibility and cost, including base load, flexible-load baseline, PV availability, price, and any outage-enabling indicators.
\end{assumption}

\begin{assumption}[Reliability semantics]
The binary variable \(I_{t,d}\) indicates whether period \((t,d)\) is in an admissible faulted state for involuntary load curtailment. In normal operation, \(I_{t,d}=0\); if reliability states are modeled endogenously, they must be linked to exogenous outage-enabling parameters and cannot serve as free feasibility switches.
\end{assumption}

\begin{assumption}[Flexible-load semantics]
The total demand is decomposed into a non-shiftable base part and a flexible part. Flexible demand may be temporally shifted within a day, but its daily energy requirement must be preserved unless an explicit interruptible-load product is introduced.
\end{assumption}

\begin{assumption}[Annual carbon semantics]
Annual carbon control is imposed on the cumulative annual grid-import emissions. It is therefore an annual aggregate condition and not, in general, a collection of independent daily caps.
\end{assumption}

\subsection{Two-Stage Stochastic Optimization Model}

\subsubsection{Planning Problem Statement}

Let \(Q(x,\xi_d)\) denote the optimal day-\(d\) operating value induced by first-stage decision \(x\), and let
\begin{align}
G(x,\xi_d) := \kappa \sum_{t\in\mathcal{T}} p^{\mathrm{buy}}_{t,d}(x,\xi_d)\,\Delta t
\label{eq:10a-carbon-functional}
\end{align}
denote the day-wise carbon-emission functional, where \(\kappa\) is the grid emission factor in kg/kWh. Only imports contribute to annual carbon accounting. Let
\begin{align}
H(x,\xi_d) := \sum_{t\in\mathcal{T}} \nu_t I_{t,d}(x,\xi_d)
\label{eq:10a-reliability-functional}
\end{align}
denote a day-wise reliability surrogate, where \(\nu_t\ge 0\) are time weights satisfying \(\sum_t \nu_t=1\).

The full-year two-stage stochastic program is
\begin{align}
\min_{x} \quad & C^{\mathrm{inv}}(x) + \sum_{d\in\mathcal{D}} \pi_d Q(x,\xi_d) + \rho_{\mathrm{CO_2}}\zeta^{\mathrm{CO_2}} + \rho_{\mathrm{rel}}\zeta^{\mathrm{rel}}
\label{eq:10a-two-stage-master-obj}\\
\text{s.t.}\quad & x \in \mathcal{X},
\label{eq:10a-first-stage-set}\\
& \sum_{d\in\mathcal{D}} \pi_d G(x,\xi_d) \le \bar E^{\mathrm{ann}}/365 + \zeta^{\mathrm{CO_2}},
\label{eq:10a-first-stage-carbon}\\
& \sum_{d\in\mathcal{D}} \pi_d H(x,\xi_d) \le R_{\mathrm{thr}} + \zeta^{\mathrm{rel}},
\label{eq:10a-first-stage-reliability}\\
& \zeta^{\mathrm{CO_2}} \ge 0,\qquad \zeta^{\mathrm{rel}} \ge 0.
\label{eq:10a-first-stage-slacks}
\end{align}

Here \(\bar E^{\mathrm{ann}}\) is the annual carbon budget and \(R_{\mathrm{thr}}\) is the admissible annual reliability threshold. The slack penalties satisfy \(\rho_{\mathrm{CO_2}},\rho_{\mathrm{rel}}>0\). For hard-constraint formulations, the two slack variables are suppressed.

\subsubsection{First-Stage Structure and Annual Couplings}

The first-stage objective \eqref{eq:10a-two-stage-master-obj} contains investment cost and expected operating cost. The coupling constraints \eqref{eq:10a-first-stage-carbon}--\eqref{eq:10a-first-stage-reliability} are annual constraints and therefore must remain outside any naive day-by-day decomposition.

The carbon constraint \eqref{eq:10a-first-stage-carbon} controls cumulative annual emissions through expected daily carbon. The reliability constraint \eqref{eq:10a-first-stage-reliability} controls the expected fraction of time periods that enter an admissible faulted state. These two quantities are structurally different, but both are scenario-coupling functionals of the second-stage recourse.

For capital-cost modeling, however, one additional consistency point is essential: the recourse term in \eqref{eq:10a-two-stage-master-obj} is an annual operating cost, so the first-stage investment term must be put on the same annual basis. Raw upfront CAPEX cannot be added directly to one-year operating cost, and dividing CAPEX only by a discount-rate scalar is not a valid net-present-value conversion. The correct comparison basis is either full-life NPV or its equivalent annual cost.

Let asset type \(a\in\{\mathrm{pv},\mathrm{bp},\mathrm{be},\mathrm{grid}\}\) denote PV, battery power, battery energy, and grid interconnection, respectively. Let \(K_a^0\) be the upfront unit investment cost, \(n_a\) the economic life in years, and \(r\ge 0\) the discount rate. Define the capital-recovery factor
\begin{align}
\operatorname{crf}(r,n) :=
\begin{cases}
\dfrac{r(1+r)^n}{(1+r)^n-1}, & r>0,\\
\dfrac{1}{n}, & r=0.
\end{cases}
\label{eq:10a-crf}
\end{align}
Then the annualized unit cost of asset \(a\) is
\begin{align}
\widehat K_a := \operatorname{crf}(r,n_a) K_a^0.
\label{eq:10a-annualized-unit-capex}
\end{align}

If the planning model uses one representative standard year and interprets the operating term as expected annual operating cost, then the consistent first-stage investment term is
\begin{align}
C^{\mathrm{inv}}(x)
= \widehat K_{\mathrm{pv}} P^{\mathrm{pv}}
+ \widehat K_{\mathrm{bp}} P^{\mathrm{bess}}
+ \widehat K_{\mathrm{be}} E^{\mathrm{bess}}
+ \widehat K_{\mathrm{grid}} P^{\mathrm{grid}}.
\label{eq:10a-annualized-investment-cost}
\end{align}

This annualized objective is equivalent to a stationary-life-cycle NPV model. Indeed, if one year of expected operation repeats over the asset life and there is no explicit salvage term, then minimizing
\begin{align}
\mathrm{NPV}(x)
= K^0(x) + \sum_{\tau=1}^{N} \frac{\mathbb{E}[Q^{\mathrm{ann}}(x,\xi)]}{(1+r)^\tau}
\label{eq:10a-npv-life-cycle}
\end{align}
is equivalent to minimizing
\begin{align}
\operatorname{crf}(r,N) K^0(x) + \mathbb{E}[Q^{\mathrm{ann}}(x,\xi)].
\label{eq:10a-eac-equivalent}
\end{align}
Therefore the annual planning objective must combine annualized investment cost with annual operating cost, not upfront investment cost with annual operating cost.

\subsubsection{Second-Stage Dispatch and Operational Constraints}

For fixed \((x,\xi_d)\), the second-stage problem is
\begin{align}
Q(x,\xi_d) := \min_{y_d} \quad & \sum_{t\in\mathcal{T}} \Big(c^{\mathrm{buy}}_{t,d} p^{\mathrm{buy}}_{t,d} - c^{\mathrm{sell}}_{t,d} p^{\mathrm{sell}}_{t,d} + c^{\mathrm{ls}} p^{\mathrm{ls}}_{t,d} + c^{\mathrm{flex}}(u^+_{t,d}+u^-_{t,d}) \\
&\qquad\\ + c^{\mathrm{thru}} p^{\mathrm{ch}}_{t,d} + c^{\mathrm{thru}} p^{\mathrm{dis}}_{t,d} + c^{\mathrm{soc}}_1 z^{(1)}_{t,d} + c^{\mathrm{soc}}_2 z^{(2)}_{t,d}\Big)
\label{eq:10a-second-stage-obj}
\end{align}
subject to the following constraints for all \(t\in\mathcal{T}\):

\paragraph{Power balance with flexible load.}
\begin{align}
p^{\mathrm{buy}}_{t,d} - p^{\mathrm{sell}}_{t,d} + p^{\mathrm{pv}}_{t,d} + p^{\mathrm{dis}}_{t,d} - p^{\mathrm{ch}}_{t,d} + p^{\mathrm{ls}}_{t,d}
= L^{\mathrm{base}}_{t,d} + \ell^{\mathrm{flex}}_{t,d}.
\label{eq:10a-power-balance}
\end{align}
Equation \eqref{eq:10a-power-balance} is the operational balance equation in which flexible load appears explicitly on the demand side. This is the correct location for flexible demand in the recourse model.

\paragraph{PV and grid limits.}
\begin{align}
0 \le p^{\mathrm{pv}}_{t,d} &\le \alpha^{\mathrm{pv}}_{t,d} P^{\mathrm{pv}},
\label{eq:10a-pv-limit}\\
0 \le p^{\mathrm{buy}}_{t,d} &\le P^{\mathrm{grid}},
\label{eq:10a-grid-import-limit}\\
0 \le p^{\mathrm{sell}}_{t,d} &\le P^{\mathrm{grid}},
\label{eq:10a-grid-export-limit}\\
p^{\mathrm{buy}}_{t,d} + p^{\mathrm{sell}}_{t,d} &\le P^{\mathrm{grid}}.
\label{eq:10a-grid-trade-hull}
\end{align}
Constraint \eqref{eq:10a-grid-trade-hull} is a linear convex-hull strengthening of the feasible grid-trade set under a common interconnection-capacity limit.

\paragraph{Battery power and energy constraints.}
\begin{align}
0 \le p^{\mathrm{ch}}_{t,d} &\le P^{\mathrm{bess}},
\label{eq:10a-battery-charge-limit}\\
0 \le p^{\mathrm{dis}}_{t,d} &\le P^{\mathrm{bess}},
\label{eq:10a-battery-discharge-limit}\\
0 \le e_{t,d} &\le E^{\mathrm{bess}},
\label{eq:10a-battery-energy-limit}\\
e_{t,d} &= e_{t-1,d} + \eta^{\mathrm{ch}} p^{\mathrm{ch}}_{t,d}\Delta t - \frac{1}{\eta^{\mathrm{dis}}} p^{\mathrm{dis}}_{t,d}\Delta t.
\label{eq:10a-battery-dynamics}
\end{align}
A small positive coefficient \(c^{\mathrm{thru}}>0\) is used in \eqref{eq:10a-second-stage-obj} as a battery-throughput regularization term representing a simplified LCOE / degradation cost.
In addition, the formal model may optionally include a low-SOC piecewise degradation regularizer through auxiliary deficit variables \(z^{(1)}_{t,d}, z^{(2)}_{t,d} \ge 0\) satisfying
\begin{align}
z^{(1)}_{t,d} &\ge \theta_1 E^{\mathrm{bess}} - e_{t,d},\\
z^{(2)}_{t,d} &\ge \theta_2 E^{\mathrm{bess}} - e_{t,d}, \qquad 0 \le \theta_1 \le \theta_2 \le 1,
\label{eq:10a-piecewise-soc-epigraph}
\end{align}
with additional objective contribution \(c^{\mathrm{soc}}_1 z^{(1)}_{t,d} + c^{\mathrm{soc}}_2 z^{(2)}_{t,d}\), where \(0 \le c^{\mathrm{soc}}_1 \le c^{\mathrm{soc}}_2\). This term remains linear and is switchable: when disabled, these auxiliary variables and coefficients are omitted from the recourse model.

For the battery, the costing basis must be unified more carefully than for PV or grid capacity, because battery economics can be charged both in first stage and in recourse. If one writes a full battery energy-capacity CAPEX term in \eqref{eq:10a-annualized-investment-cost} and also adds throughput- or SOC-based wear costs in recourse, then the cell degradation component can easily be counted twice.

To avoid this, decompose the battery energy-capacity cost into a non-wear part and a wear-related part:
\begin{align}
K_{\mathrm{be}}^0 = K_{\mathrm{be,fix}}^0 + K_{\mathrm{be,wear}}^0.
\label{eq:10a-battery-capex-split}
\end{align}
The term \(K_{\mathrm{be,fix}}^0\) covers annualized fixed ownership cost not directly caused by dispatch-induced cycling, such as enclosure, installation, balance-of-plant, and any calendar component one wishes to recover independently of operation. The term \(K_{\mathrm{be,wear}}^0\) is the portion of battery value that is consumed by cycle- and SOC-dependent degradation.

If dispatch-induced degradation is monetized explicitly in recourse, then the annualized first-stage battery term must use only the fixed component:
\begin{align}
C^{\mathrm{inv}}_{\mathrm{bess}}(x)
= \widehat K_{\mathrm{bp}} P^{\mathrm{bess}} + \widehat K_{\mathrm{be,fix}} E^{\mathrm{bess}},
\label{eq:10a-battery-fixed-investment-only}
\end{align}
and the wear-related cost is moved to recourse through
\begin{align}
D_d(y_d;E^{\mathrm{bess}})
:= c^{\mathrm{thru}} \sum_{t\in\mathcal{T}} \big(p^{\mathrm{ch}}_{t,d}+p^{\mathrm{dis}}_{t,d}\big)\Delta t
+ \sum_{t\in\mathcal{T}} \Big(c^{\mathrm{soc}}_1 z^{(1)}_{t,d} + c^{\mathrm{soc}}_2 z^{(2)}_{t,d}\Big)\Delta t.
\label{eq:10a-battery-operational-degradation-cost}
\end{align}
In this interpretation, \(c^{\mathrm{thru}}\) is a baseline cycling-wear coefficient, while the piecewise-SOC term is only an incremental low-SOC stress premium. The SOC-band term must therefore be interpreted as additional degradation beyond the baseline throughput effect, not as a second full replacement charge.

Conversely, if the model keeps the full annualized battery energy CAPEX
\begin{align}
\widehat K_{\mathrm{be}} E^{\mathrm{bess}}
\label{eq:10a-full-battery-capex-annualized}
\end{align}
in the first stage, then the recourse coefficients \(c^{\mathrm{thru}},c^{\mathrm{soc}}_1,c^{\mathrm{soc}}_2\) cannot simultaneously be interpreted as economic battery-degradation charges. In that case they may still be used as small regularization terms, but they should be documented as non-economic tie-breakers rather than cost components of the social objective.

This leads to the following no-double-counting principle.

\begin{proposition}[No-double-counting condition for battery costs]
Let the battery objective contribution be written as the sum of a first-stage annualized capacity term and a second-stage dispatch-dependent degradation term. Then a consistent economic model must satisfy exactly one of the following two regimes:
\begin{enumerate}
  \item \textbf{CAPEX-only regime:} the first stage contains the full annualized battery energy-capacity cost \(\widehat K_{\mathrm{be}} E^{\mathrm{bess}}\), and the recourse coefficients \(c^{\mathrm{thru}},c^{\mathrm{soc}}_1,c^{\mathrm{soc}}_2\) are set to zero or treated as purely numerical regularizers without economic interpretation;
  \item \textbf{Split-cost regime:} the first stage contains only \(\widehat K_{\mathrm{be,fix}} E^{\mathrm{bess}}\), while the wear-related portion \(K_{\mathrm{be,wear}}^0\) is recovered through the dispatch-dependent degradation term \eqref{eq:10a-battery-operational-degradation-cost}; in this case the SOC-band term must represent only incremental low-SOC wear beyond the baseline throughput term.
\end{enumerate}
If both the full battery energy-capacity CAPEX and a monetized throughput/SOC wear term are included simultaneously, then the battery wear cost is double counted.
\end{proposition}

\begin{proof}
The battery wear component is an economic claim on the same physical quantity: the consumable life of the electrochemical cell. Charging the full cell-replacement value in first stage and then charging a second monetized wear term in recourse amounts to recovering that same consumable life twice. The only consistent alternatives are therefore: either recover the cell value entirely through annualized ownership cost and omit economic wear charges from recourse, or split the cell cost into fixed and wear-related parts and recover the wear-related part only through dispatch-dependent degradation charges.
\end{proof}

For practical calibration, one may map a cycle-life specification into a baseline throughput-wear coefficient. If \(N^{\mathrm{efc}}\) denotes the warranted number of equivalent full cycles and \(\delta^{\mathrm{DoD}}:=\overline s-\underline s\) denotes usable depth of discharge, then one engineering approximation for the baseline marginal wear cost per kWh of counted charge-plus-discharge throughput is
\begin{align}
c^{\mathrm{thru}} \approx \frac{K_{\mathrm{be,wear}}^0}{2 N^{\mathrm{efc}} \delta^{\mathrm{DoD}}}.
\label{eq:10a-throughput-wear-calibration}
\end{align}
Under the split-cost regime, the piecewise-SOC coefficients \(c^{\mathrm{soc}}_1,c^{\mathrm{soc}}_2\) should then be calibrated only as incremental low-SOC stress premia beyond \eqref{eq:10a-throughput-wear-calibration}, not as independent full degradation surrogates.

In the implemented split-cost reruns, this calibration is automated directly in the planning builder. The economic input specifies a warranted cycle life, a reference annual cycling duty, a reference low-SOC dwell duration, and two low-SOC stress multipliers. With the current settings
\begin{align}
N^{\mathrm{efc}} &= 6000,
&n_{\mathrm{ref}}^{\mathrm{cyc}} &= 365~\text{cycles/year},
&\delta^{\mathrm{DoD}} &= 0.8,
&T_{\mathrm{ref}}^{\mathrm{soc}} &= 4~\text{h},
\end{align}
the 10-year battery energy cost implies an automatically calibrated wear share
\begin{align}
\alpha_{\mathrm{wear}} = \min\left\{1,\frac{10\times 365}{6000}\right\} = 0.6083.
\end{align}
Using the annualized battery-energy CAPEX \(\widehat K_{\mathrm{be}}\), the implemented coefficients are then
\begin{align}
c^{\mathrm{thru}} &= \frac{\alpha_{\mathrm{wear}}\widehat K_{\mathrm{be}}}{2 n_{\mathrm{ref}}^{\mathrm{cyc}} \delta^{\mathrm{DoD}}}
\approx 0.2135,\\
c^{\mathrm{soc}}_1 &= \frac{0.15}{T_{\mathrm{ref}}^{\mathrm{soc}}} c^{\mathrm{thru}} \approx 8.00\times 10^{-3},\\
c^{\mathrm{soc}}_2 &= \frac{0.40}{T_{\mathrm{ref}}^{\mathrm{soc}}} c^{\mathrm{thru}} \approx 2.13\times 10^{-2}.
\end{align}
Hence the first stage now recovers only the remaining fixed share \((1-\alpha_{\mathrm{wear}})\widehat K_{\mathrm{be}}\) of battery-energy ownership cost, while the operational replay model recovers the calibrated cycling- and low-SOC-related portion through throughput and SOC-stress charges.

\paragraph{Flexible-load realization and daily energy conservation.}
\begin{align}
\ell^{\mathrm{flex}}_{t,d} &= L^{\mathrm{flex},0}_{t,d} + u^+_{t,d} - u^-_{t,d},
\label{eq:10a-flex-decomposition}\\
\underline L^{\mathrm{flex}}_{t,d} \le \ell^{\mathrm{flex}}_{t,d} &\le \overline L^{\mathrm{flex}}_{t,d},
\label{eq:10a-flex-bounds}\\
u^+_{t,d},u^-_{t,d} &\ge 0,
\label{eq:10a-flex-modulation-nonnegative}\\
\sum_{t\in\mathcal{T}} \ell^{\mathrm{flex}}_{t,d}\Delta t &= \sum_{t\in\mathcal{T}} L^{\mathrm{flex},0}_{t,d}\Delta t,
\label{eq:10a-flex-energy-balance}\\
\sum_{t\in\mathcal{T}} u^+_{t,d}\Delta t &= \sum_{t\in\mathcal{T}} u^-_{t,d}\Delta t.
\label{eq:10a-flex-balanced-modulation}
\end{align}

\paragraph{Load shedding and reliability-state gating.}
\begin{align}
0 \le p^{\mathrm{ls}}_{t,d} &\le \big(L^{\mathrm{base}}_{t,d}+\ell^{\mathrm{flex}}_{t,d}\big) I_{t,d},
\label{eq:10a-load-shedding-gating}\\
I_{t,d} &\le Z_{t,d},
\label{eq:10a-reliability-link}\\
I_{t,d} &\in \{0,1\}.
\label{eq:10a-reliability-binary}
\end{align}
Here \(Z_{t,d}\) is an exogenous outage-enabling parameter. Under normal operation, \(Z_{t,d}=0\), hence \(I_{t,d}=0\) and \(p^{\mathrm{ls}}_{t,d}=0\). Thus involuntary load shedding is excluded in normal states.

\subsection{Carbon- and Reliability-Oriented Optimal Scenario Aggregation and Error Analysis}

\subsubsection{LP Recourse Regime}

For theoretical analysis of scenario reduction, the recourse problem is considered in the abstract form
\begin{align}
Q(x,\xi) := \min_{y} \{q(\xi)^\top y : A(\xi)y \ge b(\xi)-B(\xi)x,\ y\in\mathcal{Y}(\xi)\}.
\label{eq:10a-abstract-recourse}
\end{align}

The conditions below are sufficient, but not necessary, for the second-stage problem to be a linear program.

\begin{assumption}[Sufficient conditions for LP recourse]
The recourse problem \eqref{eq:10a-abstract-recourse} is a linear program if:
\begin{enumerate}
  \item the second-stage objective is linear in \(y\);
  \item all operational constraints are linear in \((x,y)\);
  \item the recourse set contains no binary or integer variables;
  \item the reliability state \(I_{t,d}\) is fixed by scenario data or replaced by an exogenous parameter;
  \item any battery charge/discharge exclusivity is relaxed or convexified so that the recourse region remains polyhedral;
  \item grid import/export interaction is represented by linear bounds together with the hull constraint \eqref{eq:10a-grid-trade-hull}, or by an equivalent linear polyhedral reformulation;
  \item any low-SOC degradation cost is introduced only through continuous epigraph variables such as \eqref{eq:10a-piecewise-soc-epigraph}, so no additional fixed-charge or binary logic is added.
\end{enumerate}
\end{assumption}

Assumption 5 is intentionally stronger than necessary. In particular, the second stage may still be LP after partial reformulation or conditional fixing of some logical variables. What matters is not the original modeling language, but whether the realized recourse feasible set for fixed \((x,\xi)\) is polyhedral and described without integrality restrictions.

\begin{proposition}[Parameter regime under which LP recourse is guaranteed]
For the two-stage microgrid planning model introduced above, the second-stage problem is guaranteed to be a linear program whenever the following parameter regime holds:
\begin{enumerate}
  \item \(Z_{t,d}\) is exogenous for all \((t,d)\), and \(I_{t,d}\) is either fixed to \(Z_{t,d}\) or replaced by an exogenous fault-state parameter;
  \item battery charge/discharge complementarity is not enforced by binary variables;
  \item all import prices, export remunerations, emission factors, flexible-load bounds, and PV availability factors enter linearly and are exogenous data;
  \item no block tariffs, no unit-commitment logic, and no piecewise fixed-charge operating costs appear in recourse; continuous piecewise-affine SOC penalties written through \eqref{eq:10a-piecewise-soc-epigraph} are allowed;
  \item the load-shedding admissibility condition \eqref{eq:10a-load-shedding-gating} is linear because the gating variable is exogenous;
  \item grid exchange is bounded through \eqref{eq:10a-grid-import-limit}--\eqref{eq:10a-grid-trade-hull}, which is linear;
  \item the daily storage reset is enforced by the terminal condition implicit in \eqref{eq:10a-battery-dynamics}, so for fixed first-stage capacities the full-year replay decomposes into 365 independent day problems.
\end{enumerate}
Under these conditions, for every fixed first-stage vector \(x\) and day scenario \(\xi_d\), the recourse feasible region is a polyhedron and the recourse value function is an LP value function.
\end{proposition}

\begin{proof}
Under items 1--7, all second-stage constraints \eqref{eq:10a-power-balance}--\eqref{eq:10a-reliability-binary} either remain linear equalities/inequalities in the continuous recourse variables or reduce to linear bounds with exogenous coefficients. No integrality condition remains in recourse, and the objective \eqref{eq:10a-second-stage-obj} is linear. Therefore the recourse problem is a linear program.
\end{proof}

An important weaker statement is also useful.

\begin{remark}[Conditional LP recourse]
Even if the original full model contains binary reliability states or battery mode logic, the second stage becomes LP conditionally on a fixed logical pattern. Hence LP-based sensitivity analysis can still be used locally, for example within a sequential scheme that freezes the logical state pattern or treats it as exogenous during scenario-set refinement.
\end{remark}

\begin{proposition}[Strict price spread excludes simultaneous import and export]
Consider one period \((t,d)\) of the LP recourse model with balance equation \eqref{eq:10a-power-balance}, bounds \eqref{eq:10a-grid-import-limit}--\eqref{eq:10a-grid-trade-hull}, and objective contribution
\begin{align}
c^{\mathrm{buy}}_{t,d} p^{\mathrm{buy}}_{t,d} - c^{\mathrm{sell}}_{t,d} p^{\mathrm{sell}}_{t,d}.
\end{align}
If the price spread satisfies
\begin{align}
c^{\mathrm{buy}}_{t,d} > \eta^{\mathrm{grid}} c^{\mathrm{sell}}_{t,d} \ge 0,
\label{eq:10a-strict-grid-spread}
\end{align}
then any optimal LP solution can be chosen so that at least one of \(p^{\mathrm{buy}}_{t,d}\) and \(p^{\mathrm{sell}}_{t,d}\) is zero.
\end{proposition}

\begin{proof}
Write \(\eta^{\mathrm{grid}}\) for the import efficiency coefficient appearing in \eqref{eq:10a-power-balance}. Suppose \(p^{\mathrm{buy}}_{t,d}>0\) and \(p^{\mathrm{sell}}_{t,d}>0\), and define the net grid contribution
\begin{align}
n := \eta^{\mathrm{grid}} p^{\mathrm{buy}}_{t,d} - p^{\mathrm{sell}}_{t,d}.
\end{align}
Then the replacement \((p^{\mathrm{buy}}_{t,d},p^{\mathrm{sell}}_{t,d}) \mapsto (n/\eta^{\mathrm{grid}},0)\) preserves \eqref{eq:10a-power-balance}. It is feasible for \eqref{eq:10a-grid-import-limit}--\eqref{eq:10a-grid-trade-hull} because
\begin{align}
\frac{n}{\eta^{\mathrm{grid}}} = p^{\mathrm{buy}}_{t,d} - \frac{p^{\mathrm{sell}}_{t,d}}{\eta^{\mathrm{grid}}} \le p^{\mathrm{buy}}_{t,d} \le P^{\mathrm{grid}}.
\end{align}
The objective change under this replacement is
\begin{align}
\frac{c^{\mathrm{buy}}_{t,d}}{\eta^{\mathrm{grid}}} n - \big(c^{\mathrm{buy}}_{t,d} p^{\mathrm{buy}}_{t,d} - c^{\mathrm{sell}}_{t,d} p^{\mathrm{sell}}_{t,d}\big)
= p^{\mathrm{sell}}_{t,d}\Big(c^{\mathrm{sell}}_{t,d} - \frac{c^{\mathrm{buy}}_{t,d}}{\eta^{\mathrm{grid}}}\Big) < 0,
\end{align}
This contradicts optimality. Hence simultaneous import/export cannot be optimal under \eqref{eq:10a-strict-grid-spread}.
\end{proof}

\begin{remark}[Throughput penalty as LP regularization]
A small positive throughput penalty \(c^{\mathrm{thru}}>0\) removes zero-cost battery cycling and therefore acts as a physically motivated tie-breaker among LP-optimal dispatches. However, by itself it does not replace charge/discharge complementarity: if simultaneous charging and discharging help satisfy other linear constraints, the LP optimum may still use both.
\end{remark}

\begin{proposition}[Storage-mode perturbation with SOC preservation]
Fix a first-stage vector \(x\) and a day scenario \(\xi_d\). Consider the base LP recourse model with battery dynamics \eqref{eq:10a-battery-dynamics}, bounds \eqref{eq:10a-battery-charge-limit}--\eqref{eq:10a-battery-energy-limit}, daily storage reset, and throughput penalty \(c^{\mathrm{thru}}>0\). Let \(y^{\mathrm{LP}}\) be any feasible solution that already satisfies the strict-spread conclusion of the grid-trade proposition, so simultaneous import/export does not occur. For each period, define the realized SOC increment
\begin{align}
s_{t,d} := \eta^{\mathrm{ch}} p^{\mathrm{ch}}_{t,d} - \frac{1}{\eta^{\mathrm{dis}}} p^{\mathrm{dis}}_{t,d}.
\end{align}
Define a perturbed battery pair by
\begin{align}
\widetilde p^{\mathrm{ch}}_{t,d} := \frac{(s_{t,d})_+}{\eta^{\mathrm{ch}}},
\qquad
\widetilde p^{\mathrm{dis}}_{t,d} := \eta^{\mathrm{dis}}(-s_{t,d})_+.
\label{eq:10a-storage-mode-perturbation}
\end{align}
Then the following hold:
\begin{enumerate}
  \item \(\widetilde p^{\mathrm{ch}}_{t,d}\widetilde p^{\mathrm{dis}}_{t,d}=0\) for every \(t\), so the battery is mode-feasible period by period;
  \item \(\eta^{\mathrm{ch}}\widetilde p^{\mathrm{ch}}_{t,d} - \widetilde p^{\mathrm{dis}}_{t,d}/\eta^{\mathrm{dis}} = s_{t,d}\), hence the entire SOC trajectory \(e_{t,d}\) and the daily reset condition are preserved exactly;
  \item \(\widetilde p^{\mathrm{ch}}_{t,d}\le P^{\mathrm{bess}}\) and \(\widetilde p^{\mathrm{dis}}_{t,d}\le P^{\mathrm{bess}}\) whenever the original pair satisfies \eqref{eq:10a-battery-charge-limit}--\eqref{eq:10a-battery-discharge-limit};
  \item the throughput contribution weakly decreases:
  \begin{align}
  c^{\mathrm{thru}}\big(\widetilde p^{\mathrm{ch}}_{t,d}+\widetilde p^{\mathrm{dis}}_{t,d}\big)
  \le c^{\mathrm{thru}}\big(p^{\mathrm{ch}}_{t,d}+p^{\mathrm{dis}}_{t,d}\big),
  \label{eq:10a-storage-mode-throughput-improvement}
  \end{align}
  with strict inequality whenever \(p^{\mathrm{ch}}_{t,d}>0\) and \(p^{\mathrm{dis}}_{t,d}>0\).
\end{enumerate}
Moreover, define the period-wise excess supply created by the perturbation as
\begin{equation}
\delta_{t,d} := \big(\widetilde p^{\mathrm{dis}}_{t,d}-\widetilde p^{\mathrm{ch}}_{t,d}\big) - \big(p^{\mathrm{dis}}_{t,d}-p^{\mathrm{ch}}_{t,d}\big) \ge 0.
\label{eq:10a-storage-mode-delta}
\end{equation}
If there exist nonnegative compensation variables \((\Delta^{\mathrm{buy}}_{t,d},\Delta^{\mathrm{sell}}_{t,d},\Delta^{\mathrm{pv}}_{t,d})\) such that
\begin{equation}
\begin{aligned}
\Delta^{\mathrm{buy}}_{t,d} + \Delta^{\mathrm{sell}}_{t,d} + \Delta^{\mathrm{pv}}_{t,d} &= \delta_{t,d}, \\
0 \le \Delta^{\mathrm{buy}}_{t,d} &\le p^{\mathrm{buy}}_{t,d}, \\
0 \le \Delta^{\mathrm{sell}}_{t,d} &\le P^{\mathrm{grid}} - p^{\mathrm{sell}}_{t,d}, \\
0 \le \Delta^{\mathrm{pv}}_{t,d} &\le p^{\mathrm{pv}}_{t,d},
\end{aligned}
\label{eq:10a-storage-mode-compensation-balance}
\end{equation}
and the resulting adjusted grid/PV variables do not create simultaneous import and export, then the full perturbed solution remains feasible and its objective value does not increase. In particular, this condition is satisfied whenever the perturbation can be absorbed by some combination of reduced import, increased export, or additional PV curtailment.
\end{proposition}

\begin{proof}
By construction, \eqref{eq:10a-storage-mode-perturbation} gives \(\widetilde p^{\mathrm{ch}}_{t,d}\widetilde p^{\mathrm{dis}}_{t,d}=0\) and
\begin{align}
\eta^{\mathrm{ch}}\widetilde p^{\mathrm{ch}}_{t,d} - \frac{1}{\eta^{\mathrm{dis}}}\widetilde p^{\mathrm{dis}}_{t,d}
= (s_{t,d})_+ - (-s_{t,d})_+
= s_{t,d}.
\end{align}
Hence the battery dynamics \eqref{eq:10a-battery-dynamics} and therefore the entire SOC trajectory are preserved exactly. Feasibility of the power bounds follows from
\begin{align}
0 \le \widetilde p^{\mathrm{ch}}_{t,d} \le p^{\mathrm{ch}}_{t,d} \le P^{\mathrm{bess}},
\qquad
0 \le \widetilde p^{\mathrm{dis}}_{t,d} \le p^{\mathrm{dis}}_{t,d} \le P^{\mathrm{bess}},
\end{align}
because if \(s_{t,d}\ge 0\), then
\begin{align}
\widetilde p^{\mathrm{ch}}_{t,d} = \frac{s_{t,d}}{\eta^{\mathrm{ch}}} = p^{\mathrm{ch}}_{t,d} - \frac{p^{\mathrm{dis}}_{t,d}}{\eta^{\mathrm{ch}}\eta^{\mathrm{dis}}} \le p^{\mathrm{ch}}_{t,d},
\qquad
\widetilde p^{\mathrm{dis}}_{t,d}=0,
\end{align}
and if \(s_{t,d}<0\), then
\begin{align}
\widetilde p^{\mathrm{ch}}_{t,d}=0,
\qquad
\widetilde p^{\mathrm{dis}}_{t,d} = -\eta^{\mathrm{dis}} s_{t,d} = p^{\mathrm{dis}}_{t,d} - \eta^{\mathrm{ch}}\eta^{\mathrm{dis}} p^{\mathrm{ch}}_{t,d} \le p^{\mathrm{dis}}_{t,d}.
\end{align}
The throughput inequality follows from the same case split. In particular,
\begin{align}
\widetilde p^{\mathrm{ch}}_{t,d}+\widetilde p^{\mathrm{dis}}_{t,d} \le p^{\mathrm{ch}}_{t,d}+p^{\mathrm{dis}}_{t,d},
\end{align}
with strict inequality when both components are positive.

Equation \eqref{eq:10a-storage-mode-delta} is immediate from the construction and the fact that replacing simultaneous charge/discharge by a one-sided pair with the same SOC increment can only increase the battery's net power injection into the balance equation. If the compensation variables in \eqref{eq:10a-storage-mode-compensation-balance} exist, then reducing import, increasing export, or curtailing PV by the prescribed amounts restores \eqref{eq:10a-power-balance} without affecting SOC feasibility. The grid objective weakly decreases because import is reduced and/or export is increased under the strict-spread regime, while extra PV curtailment is cost-free in \eqref{eq:10a-second-stage-obj}. Together with \eqref{eq:10a-storage-mode-throughput-improvement}, this shows that the full objective does not increase.
\end{proof}

\begin{proposition}[Epigraph form of the SOC band cost]
For fixed \((x,\xi_d)\), define the low-SOC penalty by the epigraph system \eqref{eq:10a-piecewise-soc-epigraph} with coefficients \(0 \le c^{\mathrm{soc}}_1 \le c^{\mathrm{soc}}_2\). Then, for every feasible state trajectory \(e_{t,d}\), the minimizing auxiliary variables are given by
\begin{align}
z^{(1)}_{t,d} &= \max\{\theta_1 E^{\mathrm{bess}} - e_{t,d},0\},\\
z^{(2)}_{t,d} &= \max\{\theta_2 E^{\mathrm{bess}} - e_{t,d},0\},
\label{eq:10a-piecewise-soc-minimizer}
\end{align}
and the induced penalty is a convex continuous piecewise-affine function of \(e_{t,d}\).
\end{proposition}

\begin{proof}
Because the coefficients of \(z^{(1)}_{t,d},z^{(2)}_{t,d}\) in \eqref{eq:10a-second-stage-obj} are nonnegative, any objective-minimizing solution drives each auxiliary variable to its smallest feasible value. The lower bounds in \eqref{eq:10a-piecewise-soc-epigraph} together with nonnegativity therefore imply \eqref{eq:10a-piecewise-soc-minimizer}. Since the maximum of affine functions is convex and piecewise affine, the resulting SOC penalty is also convex and piecewise affine.
\end{proof}

\begin{theorem}[Band-cost changes preserve LP/MILP equivalence]
Fix a first-stage vector \(x\) and a day scenario \(\xi_d\). Consider the recourse model defined by \eqref{eq:10a-power-balance}--\eqref{eq:10a-load-shedding-gating} under the LP-recourse parameter regime stated above, together with the strict price-spread condition \eqref{eq:10a-strict-grid-spread}, the grid-trade hull \eqref{eq:10a-grid-trade-hull}, the daily storage reset, and the throughput penalty \(c^{\mathrm{thru}}>0\). Let the corresponding MILP impose battery-mode exclusivity, while the LP uses its continuous relaxation. Assume that the balance-compensation condition of the storage-mode perturbation proposition holds period by period. Then for any thresholds \(0 \le \theta_1 \le \theta_2 \le 1\) and any nonnegative band costs \(0 \le c^{\mathrm{soc}}_1 \le c^{\mathrm{soc}}_2\), the recourse LP and the corresponding MILP are value-equivalent.
\end{theorem}

\begin{proof}
The MILP feasible set is a subset of the LP feasible set, so only the possibility of a strictly better LP optimum must be excluded. Suppose by contradiction that, after adding the SOC band penalty, there exists an LP-optimal solution \(y^{\mathrm{LP}}\) with objective value strictly smaller than every MILP-feasible solution.

First, by the strict-spread proposition, the grid-trade pair in \(y^{\mathrm{LP}}\) may be chosen so that simultaneous import and export do not occur. Next, replace the auxiliary variables in \(y^{\mathrm{LP}}\) by their minimal epigraph values from \eqref{eq:10a-piecewise-soc-minimizer}; this does not increase the LP objective. Then apply the storage-mode perturbation proposition period by period. This yields a battery-mode-feasible solution \(\bar y\) that is MILP-feasible, preserves the SOC trajectory \(e\), and restores power balance through feasible grid/PV compensation.

Because \(\bar y\) has the same state trajectory \(e\), its minimizing SOC-band auxiliaries coincide with those of \(y^{\mathrm{LP}}\) through \eqref{eq:10a-piecewise-soc-minimizer}. Hence the additional contribution
\begin{align}
\sum_t \big(c^{\mathrm{soc}}_1 z^{(1)}_{t,d} + c^{\mathrm{soc}}_2 z^{(2)}_{t,d}\big)
\end{align}
is exactly the same for \(y^{\mathrm{LP}}\) and \(\bar y\).

The strict-spread proposition and the storage-mode perturbation proposition together show that every LP-optimal solution can be converted into a MILP-feasible solution without increasing the base objective. Since the SOC-band term is unchanged by preservation of \(e\), the same conversion does not increase the augmented objective either. This yields a MILP-feasible solution with objective value no larger than the LP optimum, contradicting the assumption of a strict LP advantage. Hence the LP and MILP recourse problems are value-equivalent for every choice of nonnegative band costs and thresholds satisfying the ordering constraints.
\end{proof}

\begin{remark}[Interpretation of the band-cost theorem]
The theorem shows that changing \((\theta_1,\theta_2,c^{\mathrm{soc}}_1,c^{\mathrm{soc}}_2)\) changes the common recourse value function, but it does not by itself create a new LP--MILP gap. The band term only reweights a continuous state-dependent penalty. Consequently, once the perturbation-based LP/MILP equivalence has been established for the base model, it survives every such SOC-band reparameterization.
\end{remark}

\begin{proposition}[Convex piecewise-affine recourse]
Under Assumption 5, the value function \(Q(x,\xi)\) is convex and piecewise affine in \(x\) for each fixed \(\xi\). Moreover, if \(\lambda(x,\xi)\) is an optimal dual solution, then
\begin{align}
g(x,\xi) = -B(\xi)^\top \lambda(x,\xi)
\label{eq:10a-dual-subgradient}
\end{align}
is a subgradient of \(Q(\cdot,\xi)\) at \(x\).
\end{proposition}

\begin{proof}
For fixed \(\xi\), the right-hand side of \eqref{eq:10a-abstract-recourse} depends affinely on \(x\). By linear programming duality, the optimal value is the pointwise supremum of affine functions in \(x\), indexed by dual-feasible multipliers. Hence \(Q(\cdot,\xi)\) is convex and piecewise affine. The subgradient representation follows directly from the dual value function.
\end{proof}

\subsubsection{Planning-Functional Reduced-Measure Construction}

Let \(\widehat{\mathbb{P}}\) denote the empirical full-year distribution on the day space \(\Xi\), and let \(\Pi_S\) denote the set of probability measures supported on at most \(S\) retained scenarios. The scenario aggregation problem is formulated as
\begin{align}
\min_{\widetilde{\mathbb{P}}\in\Pi_S}\ \sup_{x\in\mathcal{X}} \Bigg(
&\alpha_Q \big|\mathbb{E}_{\widehat{\mathbb{P}}}[Q(x,\xi)]-\mathbb{E}_{\widetilde{\mathbb{P}}}[Q(x,\xi)]\big| \\
+&\alpha_G \big|\mathbb{E}_{\widehat{\mathbb{P}}}[G(x,\xi)]-\mathbb{E}_{\widetilde{\mathbb{P}}}[\widetilde G(x,\xi)]\big| \\
+&\alpha_H \big|\mathbb{E}_{\widehat{\mathbb{P}}}[H(x,\xi)]-\mathbb{E}_{\widetilde{\mathbb{P}}}[\widetilde H(x,\xi)]\big|
\Bigg),
\label{eq:10a-optimal-scenario-aggregation}
\end{align}
where \(\widetilde G\) and \(\widetilde H\) denote the carbon and reliability functionals used by the reduced model.

Problem \eqref{eq:10a-optimal-scenario-aggregation} states directly that the reduced scenario set should approximate three objects simultaneously: expected operating value, expected carbon emissions, and expected reliability burden.

Because \eqref{eq:10a-optimal-scenario-aggregation} is generally intractable, a feature map is introduced
\begin{align}
\phi(\xi_d) := \Big[
&L^{\mathrm{base}}_{1,d},\dots,L^{\mathrm{base}}_{T,d},
L^{\mathrm{flex},0}_{1,d},\dots,L^{\mathrm{flex},0}_{T,d},
\alpha^{\mathrm{pv}}_{1,d},\dots,\alpha^{\mathrm{pv}}_{T,d},
\lambda^{\mathrm{buy}}_{1,d},\dots,\lambda^{\mathrm{buy}}_{T,d},
\lambda^{\mathrm{sell}}_{1,d},\dots,\lambda^{\mathrm{sell}}_{T,d}, \\
&Z_{1,d},\dots,Z_{T,d},
\sum_t L^{\mathrm{flex},0}_{t,d}\Delta t,
\max_t \big(L^{\mathrm{base}}_{t,d}+L^{\mathrm{flex},0}_{t,d}-\alpha^{\mathrm{pv}}_{t,d}P^{\mathrm{pv,ref}}\big), \\
&\sum_t \big(L^{\mathrm{base}}_{t,d}+L^{\mathrm{flex},0}_{t,d}-\alpha^{\mathrm{pv}}_{t,d}P^{\mathrm{pv,ref}}\big)_+\Delta t
\Big],
\label{eq:10a-feature-map}
\end{align}
and a metric \(\rho(\xi,\xi') := \|\phi(\xi)-\phi(\xi')\|\). The last two components in \eqref{eq:10a-feature-map} are included to preserve peak stress and deficit-duration stress. The flexible-load baseline enters explicitly, which is necessary because the recourse balance \eqref{eq:10a-power-balance} depends on \(\ell^{\mathrm{flex}}_{t,d}\).

The tractable surrogate scenario-selection problem is then the optimal transport reduction
\begin{align}
\min_{\widetilde{\mathbb{P}}\in\Pi_S} W_1(\widehat{\mathbb{P}},\widetilde{\mathbb{P}};\rho),
\label{eq:10a-transport-reduction}
\end{align}
where \(W_1\) is the 1-Wasserstein distance induced by \(\rho\).

Problem \eqref{eq:10a-transport-reduction} is substantially easier than \eqref{eq:10a-optimal-scenario-aggregation}, but it can still be computationally demanding when the day-feature dimension is high or when the candidate support set is large. This motivates an iterative support-expansion strategy in which one does not attempt to solve the full reduction problem to global optimality in one shot.

Let \(\mathbb{P}_k\) denote the reduced scenario distribution at iteration \(k\), and let \(x_k\) be the corresponding planning solution. A practical sequential alternative is:
\begin{enumerate}
  \item solve the reduced planning problem under \(\mathbb{P}_k\) and obtain \(x_k\);
  \item monitor the statistical distance between the current reduced support and the full-year empirical measure, for example through probability-weighted feature-space distances induced by \eqref{eq:10a-feature-map};
  \item replay \(x_k\) on the full-year scenario set or on a large validation pool, using independent day subproblems whenever the storage state is reset daily;
  \item identify the scenarios that create the largest carbon, reliability, or recourse mismatch;
  \item augment the support of \(\mathbb{P}_k\) with those scenarios to obtain \(\mathbb{P}_{k+1}\);
  \item stop if replay becomes feasible, or if the planning solution stabilizes in the sense that \(\|x_{k+1}-x_k\|\le \varepsilon_x\), while continuing to log the support-distance statistics across iterations.
\end{enumerate}

This sequential mechanism is often preferable computationally because it transforms one large reduction problem into a sequence of smaller planning solves and targeted support enrichments.

\subsubsection{Approximation Error Bounds}

The error analysis for operating value, carbon accounting, and reliability approximation is now formalized.

\begin{assumption}[Uniform stability]
There exist constants \(L_Q,L_G,L_H<\infty\) such that, for all \(x\in\mathcal{X}\) and \(\xi,\xi'\in\Xi\),
\begin{align}
|Q(x,\xi)-Q(x,\xi')| &\le L_Q\,\rho(\xi,\xi'),
\label{eq:10a-lipschitz-q}\\
|G(x,\xi)-G(x,\xi')| &\le L_G\,\rho(\xi,\xi'),
\label{eq:10a-lipschitz-g}\\
|H(x,\xi)-H(x,\xi')| &\le L_H\,\rho(\xi,\xi').
\label{eq:10a-lipschitz-h}
\end{align}
Moreover, suppose the reduced model satisfies
\begin{align}
|G(x,\xi)-\widetilde G(x,\xi)| &\le \varepsilon_G^{\mathrm{surr}},
\label{eq:10a-carbon-surrogate-error}\\
|H(x,\xi)-\widetilde H(x,\xi)| &\le \varepsilon_H^{\mathrm{surr}}.
\label{eq:10a-reliability-surrogate-error}
\end{align}
\end{assumption}

\begin{theorem}[Combined aggregation error bound]
Under Assumption 6, for any reduced measure \(\widetilde{\mathbb{P}}\in\Pi_S\),
\begin{align}
\sup_{x\in\mathcal{X}} \Bigg(
&\alpha_Q \big|\mathbb{E}_{\widehat{\mathbb{P}}}[Q(x,\xi)]-\mathbb{E}_{\widetilde{\mathbb{P}}}[Q(x,\xi)]\big| \\
+&\alpha_G \big|\mathbb{E}_{\widehat{\mathbb{P}}}[G(x,\xi)]-\mathbb{E}_{\widetilde{\mathbb{P}}}[\widetilde G(x,\xi)]\big| \\
+&\alpha_H \big|\mathbb{E}_{\widehat{\mathbb{P}}}[H(x,\xi)]-\mathbb{E}_{\widetilde{\mathbb{P}}}[\widetilde H(x,\xi)]\big|
\Bigg)
&\le (\alpha_Q L_Q + \alpha_G L_G + \alpha_H L_H)
W_1(\widehat{\mathbb{P}},\widetilde{\mathbb{P}};\rho) \\
&\qquad + \alpha_G \varepsilon_G^{\mathrm{surr}} + \alpha_H \varepsilon_H^{\mathrm{surr}}.
\label{eq:10a-combined-error-bound}
\end{align}
\end{theorem}

\begin{proof}
For the recourse term, Kantorovich--Rubinstein duality and \eqref{eq:10a-lipschitz-q} yield
\begin{align}
\big|\mathbb{E}_{\widehat{\mathbb{P}}}[Q]-\mathbb{E}_{\widetilde{\mathbb{P}}}[Q]\big|
\le L_Q W_1(\widehat{\mathbb{P}},\widetilde{\mathbb{P}};\rho).
\end{align}
The same argument with \eqref{eq:10a-lipschitz-g} and \eqref{eq:10a-lipschitz-h} gives the bounds for the exact carbon and reliability functionals. For the reduced functionals, add and subtract \(\mathbb{E}_{\widetilde{\mathbb{P}}}[G]\) and \(\mathbb{E}_{\widetilde{\mathbb{P}}}[H]\), respectively, and use \eqref{eq:10a-carbon-surrogate-error}--\eqref{eq:10a-reliability-surrogate-error}. Summing the three estimates with weights \(\alpha_Q,\alpha_G,\alpha_H\) yields \eqref{eq:10a-combined-error-bound}.
\end{proof}

\begin{corollary}[Safe tightening of carbon and reliability thresholds]
Suppose the reduced model enforces
\begin{align}
\mathbb{E}_{\widetilde{\mathbb{P}}}[\widetilde G(x,\xi)] &\le \bar E^{\mathrm{ann}}/365 - \Delta_G,
\label{eq:10a-safe-carbon}\\
\mathbb{E}_{\widetilde{\mathbb{P}}}[\widetilde H(x,\xi)] &\le R_{\mathrm{thr}} - \Delta_H.
\label{eq:10a-safe-reliability}
\end{align}
If
\begin{align}
\Delta_G &\ge L_G W_1(\widehat{\mathbb{P}},\widetilde{\mathbb{P}};\rho) + \varepsilon_G^{\mathrm{surr}},
\label{eq:10a-safe-carbon-margin}\\
\Delta_H &\ge L_H W_1(\widehat{\mathbb{P}},\widetilde{\mathbb{P}};\rho) + \varepsilon_H^{\mathrm{surr}},
\label{eq:10a-safe-reliability-margin}
\end{align}
then the exact full-year functionals satisfy
\begin{align}
\mathbb{E}_{\widehat{\mathbb{P}}}[G(x,\xi)] &\le \bar E^{\mathrm{ann}}/365,
\label{eq:10a-exact-carbon-feasible}\\
\mathbb{E}_{\widehat{\mathbb{P}}}[H(x,\xi)] &\le R_{\mathrm{thr}}.
\label{eq:10a-exact-reliability-feasible}
\end{align}
\end{corollary}

\begin{proof}
Apply the carbon and reliability parts of \eqref{eq:10a-combined-error-bound} separately with \((\alpha_Q,\alpha_G,\alpha_H)=(0,1,0)\) and \((0,0,1)\), respectively, then substitute \eqref{eq:10a-safe-carbon}--\eqref{eq:10a-safe-reliability-margin}.
\end{proof}

Theorem 3 and Corollary 1 show that scenario reduction for this problem should be driven by a metric that preserves the combined sensitivity of recourse cost, carbon emissions, and reliability burden. A reduced scenario set that is accurate only in Euclidean profile distance need not be accurate in the planning functionals that actually determine the first-stage decision.

The same theorem also explains why sequential support expansion can work well in practice: if each enrichment step reduces the effective Wasserstein mismatch in the stress-relevant region of scenario space, then the value, carbon, and reliability errors all decrease together. However, this is a stability statement, not a global optimization guarantee for the sequence \(\{\mathbb{P}_k\}\).

\begin{proposition}[Finite termination of nested support expansion over a finite day pool]
Suppose the candidate scenario pool is finite, and consider an iterative procedure in which:
\begin{enumerate}
  \item the support of \(\mathbb{P}_{k+1}\) contains the support of \(\mathbb{P}_k\);
  \item at each unsuccessful iteration, at least one previously unselected scenario with positive replay violation score is added;
  \item the stopping test requires both
  \begin{align}
  \|x_{k+1}-x_k\| \le \varepsilon_x
  \label{eq:10a-x-stability-stop}
  \end{align}
  and replay residuals below prescribed tolerances.
\end{enumerate}
Then the procedure terminates in finitely many iterations, either because the stopping test is satisfied or because the full candidate pool has been exhausted.
\end{proposition}

\begin{proof}
The candidate scenario pool is finite. By construction, each unsuccessful iteration strictly enlarges the support. Therefore the algorithm cannot add new scenarios indefinitely. Hence finite termination follows.
\end{proof}

\begin{remark}[What finite termination does not imply]
The proposition above guarantees only finite termination of the support-expansion process over a finite day pool. It does not imply global optimality of the terminal plan for the original full-year stochastic program, nor global optimality of the induced sequence \(\{\mathbb{P}_k\}\) for the reduction problem \eqref{eq:10a-optimal-scenario-aggregation}. The stopping condition \eqref{eq:10a-x-stability-stop} is therefore best interpreted as a practical stationarity criterion rather than a proof of globally optimal aggregation.
\end{remark}

\begin{remark}[Why global convergence is difficult]
The mapping
\begin{align}
\mathbb{P}_k \mapsto x_k \mapsto \text{replay violations} \mapsto \mathbb{P}_{k+1}
\end{align}
is generally nonlinear and nonsmooth. It depends on the planning optimum, on the violation-selection rule, and on the support-update operator. Unless strong convexity, unique minimizers, and a carefully controlled enrichment rule are imposed, one should not expect a global convergence theorem to the optimizer of \eqref{eq:10a-optimal-scenario-aggregation}.
\end{remark}

\subsection{Sequential Scenario-Measure Expansion Algorithm}

The foregoing analysis leads to the following planning-oriented scenario aggregation algorithm.

\paragraph{Algorithm 1: Carbon- and reliability-oriented scenario reduction for planning.}

\begin{algorithmic}[1]
\State \textbf{Input:} full-year day set \(\mathcal{D}\), support size \(S\), carbon budget \(\bar E^{\mathrm{ann}}\), reliability threshold \(R_{\mathrm{thr}}\)
\State Construct day features \(\phi(\xi_d)\) according to \eqref{eq:10a-feature-map}
\State Solve the optimal transport reduction \eqref{eq:10a-transport-reduction} and obtain reduced measure \(\widetilde{\mathbb{P}}\in\Pi_S\)
\State Estimate or upper-bound \(L_G, L_H, \varepsilon_G^{\mathrm{surr}}, \varepsilon_H^{\mathrm{surr}}\)
\State Compute safe margins \(\Delta_G,\Delta_H\) from \eqref{eq:10a-safe-carbon-margin}--\eqref{eq:10a-safe-reliability-margin}
\State Solve the reduced two-stage planning model with tightened thresholds \eqref{eq:10a-safe-carbon}--\eqref{eq:10a-safe-reliability}
\State Replay the obtained first-stage plan on the full-year day set \(\mathcal{D}\)
\State Evaluate annual carbon feasibility, reliability feasibility, and normal-state load shedding
\If{either annual carbon or reliability constraints are violated in replay}
  \State Enrich the feature map or increase support size \(S\)
  \State Return to Step 3
\EndIf
\State \textbf{Output:} planning decision \(x\) and reduced scenario set \(\widetilde{\mathbb{P}}\)
\end{algorithmic}

\paragraph{Algorithm 2: Dual-guided scenario enrichment.}

When the reduced planning problem satisfies the LP recourse conditions, the convex piecewise-affine recourse proposition provides subgradients through \eqref{eq:10a-dual-subgradient}. These subgradients can be used to refine the reduced scenario set.

\begin{algorithmic}[1]
\State Start from an initial reduced set \(\widetilde{\mathbb{P}}^0\)
\For{iteration \(k=0,1,2,\dots\)}
  \State Solve the reduced planning problem and obtain \(x^k\)
  \State Compute dual subgradients \(g(x^k,\xi_s)\) for retained scenarios
  \State For omitted days, evaluate mismatch scores based on recourse value, carbon error, reliability error, and dual-sensitivity deviation
  \State Insert the highest-score days into the retained set and recluster the remaining days
  \State Terminate if the replay errors satisfy the certified margins in Corollary 1
\EndFor
\end{algorithmic}

Algorithm 1 yields a certified reduced-scenario planning procedure, while Algorithm 2 provides an adaptive refinement mechanism when additional accuracy is required. Their common principle is that scenario aggregation should be guided by planning functionals and their sensitivities, not merely by profile similarity.

\paragraph{Algorithm 3: Sequential scenario-measure expansion.}

The following variant is often more scalable than solving \eqref{eq:10a-transport-reduction} aggressively from the outset. To state it more formally, introduce the replay residuals
\begin{align}
r_G(x) &:= \Big[\mathbb{E}_{\widehat{\mathbb{P}}}[G(x,\xi)] - \bar E^{\mathrm{ann}}/365\Big]_+,
\label{eq:10a-replay-carbon-residual}\\
r_H(x) &:= \Big[\mathbb{E}_{\widehat{\mathbb{P}}}[H(x,\xi)] - R_{\mathrm{thr}}\Big]_+,
\label{eq:10a-replay-reliability-residual}\\
r_Q(x,\mathbb{P}) &:= \Big|\mathbb{E}_{\widehat{\mathbb{P}}}[Q(x,\xi)] - \mathbb{E}_{\mathbb{P}}[Q(x,\xi)]\Big|.
\label{eq:10a-replay-value-residual}
\end{align}
Define the merit function
\begin{align}
\mathcal{M}_{\Gamma}(x,\mathbb{P}) := C^{\mathrm{inv}}(x) + \mathbb{E}_{\mathbb{P}}[Q(x,\xi)] + \Gamma_G r_G(x) + \Gamma_H r_H(x) + \Gamma_Q r_Q(x,\mathbb{P}),
\label{eq:10a-merit-function}
\end{align}
where \(\Gamma := (\Gamma_G,\Gamma_H,\Gamma_Q)\) is a vector of strictly positive penalty parameters. The role of \(\mathcal{M}_{\Gamma}\) is to combine reduced-model optimality and full-year replay feasibility into a single descent target.

In the implemented Case 4 workflow, the reduced-support iterate is also accompanied by a monitored support-distance statistic
\begin{align}
d_{\mathrm{supp}}(\mathbb{P}_k,\widehat{\mathbb{P}})
:= \sum_{d\in\mathcal{D}} \pi_d \, \rho\big(\xi_d,\xi_{\mathrm{rep}(d;k)}\big),
\label{eq:10a-support-monitor-distance}
\end{align}
where \(\xi_{\mathrm{rep}(d;k)}\) denotes the representative scenario currently covering day \(d\). In practice, both this probability-weighted average distance and the corresponding maximum covered distance are recorded at every iteration in order to track how support enrichment changes the discrepancy between the planning support and the full-year empirical distribution.

\begin{algorithmic}[1]
\State Choose an initial reduced measure \(\mathbb{P}_0\) supported on a small retained scenario set
\For{iteration \(k=0,1,2,\dots\)}
  \State Solve the reduced planning problem under \(\mathbb{P}_k\) and obtain \(x_k\)
  \State Record the probability-weighted and maximum support-distance statistics for the current retained support
  \If{\(k\ge 1\) and \(\|x_k-x_{k-1}\| \le \varepsilon_x\)}
    \State Stop with planning-solution-stable
  \EndIf
  \State Replay \(x_k\) on the full-year scenario pool
  \State Compute scenario violation scores combining carbon mismatch, reliability mismatch, and recourse-cost mismatch
  \State Construct \(\mathbb{P}_{k+1}\) by adding the highest-score scenarios, reassigning every full-year day to its nearest retained representative under \(\rho\), and reweighting each support point by the total probability mass of its reassigned cluster
  \If{replay residuals are below tolerance}
    \State Stop with replay-feasible
  \EndIf
\EndFor
\end{algorithmic}

Algorithm 3 is computationally attractive because it postpones large-scale transport reduction and concentrates effort on scenarios revealed by simulation to be planning-critical. In the implemented Case 4 workflow, once new scenarios are inserted, the full-year empirical measure is reassigned to the current retained support by nearest-neighbor matching in the feature metric \(\rho\), and the reduced probabilities are recomputed from the resulting assigned masses; this is a Wasserstein-oriented reweighting step rather than a split-only local correction. The added planning-solution-stable test prevents unnecessary extra solves once the reduced planning iterate has effectively ceased to move, and the workflow does not append a default full-year polish solve after replay-feasible stopping. The support-distance monitor provides a direct diagnostic of how quickly the enriched retained support approaches the full-year empirical geometry. Its main limitation is theoretical: neither replay-feasible stopping nor planning-solution-stable stopping establishes global optimality; both remain empirical termination criteria for the sequentially enriched planning problem.

\paragraph{Algorithm 4: Merit-monitored support update.}

For the present study, a lighter formulation is more appropriate than a full descent method, because the implemented Case 4 workflow terminates after only two planning-replay cycles. The useful abstraction is therefore not a separate optimization algorithm over \((x,\mathbb{P})\), but a monitoring template in which the scalar merit \(\mathcal{M}_{\Gamma}\) is used to judge whether support enrichment is helping.

\begin{algorithmic}[1]
\State Choose an initial reduced measure \(\mathbb{P}_0\) and penalty parameters \(\Gamma_G,\Gamma_H,\Gamma_Q>0\)
\For{iteration \(k=0,1,2,\dots\)}
  \State Solve the reduced planning problem under \(\mathbb{P}_k\) and obtain \(x_k\)
  \State Replay \(x_k\) on the full-year scenario pool and evaluate \(\mathcal{M}_{\Gamma}(x_k,\mathbb{P}_k)\)
  \State Rank omitted scenarios by replay violation score
  \State Build a candidate enriched support \(\widehat{\mathbb{P}}_{k+1}\) by adding the highest-score scenarios and reweighting the retained support
  \State Resolve the reduced planning problem under \(\widehat{\mathbb{P}}_{k+1}\) and replay the resulting plan
  \If{the replay becomes feasible or the planning solution is stable}
    \State Stop
  \EndIf
\EndFor
\end{algorithmic}

Algorithm 4 should therefore be read only as a merit-monitored template. In the code path used for Case 4, the merit value is recorded as a diagnostic scalar, but there is no separate line-search, no damped re-solve loop, and no attempt to prove monotone descent over all iterations.

\paragraph{Algorithm 5: Implemented trust-style support controller.}

The implemented controller is intentionally simpler than a textbook Wasserstein trust-region method. The radius \(\Delta_k\) is not used to solve a constrained measure-update subproblem. Instead, it is interpreted as an integer-valued support-update budget controlling how many top-ranked replay-critical scenarios may be inserted at iteration \(k\). After insertion, all 365 days are reassigned to their nearest retained representatives under \(\rho\), and the representative probabilities are recomputed from the reassigned masses exactly as in Algorithm 3.

In this implemented variant, the predicted reduction is a proxy constructed from support-geometry improvement and captured replay-score mass rather than from an exact local subproblem model. Writing
\begin{align}
\psi(\mathbb{P}_k) := d_{\mathrm{supp}}(\mathbb{P}_k,\widehat{\mathbb{P}}) + 0.25\, d^{\max}_{\mathrm{supp}}(\mathbb{P}_k,\widehat{\mathbb{P}}),
\end{align}
the code forms a normalized prediction term from the relative decrease in \(\psi\) and from the fraction of total positive replay score captured by the inserted scenarios. The resulting predicted reduction is then scaled by the current merit value. The actual reduction remains
\begin{align}
\mathrm{Ared}_k := \mathcal{M}_{\Gamma}(x_k,\mathbb{P}_k)-\mathcal{M}_{\Gamma}(x_{k+1},\mathbb{P}_{k+1}),
\label{eq:10a-actual-reduction}
\end{align}
and the monitoring ratio is
\begin{align}
\rho_k := \frac{\mathrm{Ared}_k}{\mathrm{Pred}_k}.
\label{eq:10a-trust-ratio}
\end{align}

\begin{algorithmic}[1]
\State Choose an initial reduced support \(\mathbb{P}_0\), an initial insertion budget \(\Delta_0\), and thresholds \(0<\eta_1<\eta_2<1\)
\For{iteration \(k=0,1,2,\dots\)}
  \State Solve the reduced planning problem under \(\mathbb{P}_k\) and replay the resulting plan
  \State Evaluate the merit value \(\mathcal{M}_{\Gamma}(x_k,\mathbb{P}_k)\)
  \State Insert the top \(\max\{1,\mathrm{round}(\Delta_k)\}\) replay-critical scenarios into the retained support
  \State Reassign all full-year days to their nearest retained representatives and recompute support weights
  \State Estimate \(\mathrm{Pred}_k\) from support-distance contraction and captured replay-score mass
  \State At the next replayed iterate, compute \(\mathrm{Ared}_k\) and \(\rho_k\)
  \If{\(\rho_k \ge \eta_2\)}
    \State Increase the next insertion budget
  \ElsIf{\(0 \le \rho_k < \eta_1\)}
    \State Decrease the next insertion budget
  \Else
    \State Keep the insertion budget unchanged
  \EndIf
  \If{replay is feasible or the planning solution is stable}
    \State Stop
  \EndIf
\EndFor
\end{algorithmic}

Two implementation details are important. First, the support sets remain nested: scenarios inserted at one iteration are not removed later. Second, unlike a full trust-region method, a low ratio does not reject the current support update or trigger a restoration solve; it only shrinks the insertion budget used for the next enrichment step. Hence the controller is best interpreted as a lightweight safeguard against overly aggressive support growth, not as a full trust-region algorithm with a separate constrained subproblem and subsequence-convergence guarantees.

This distinction matters practically because, in the present Case 4 experiments, both the plain sequential update and the trust-enabled variant terminate after two iterations. The trust-style controller is therefore useful mainly as a diagnostic layer that records merit, predicted reduction, actual reduction, and radius adaptation in a unified way; it is not needed to explain convergence of a long iterative process in this data set.

\subsection{Experimental Design}

The numerical study is organized around four cases whose purpose is to compare not only objective values and runtimes, but also whether reduced models preserve the feasibility structure of the original annual planning problem under the corrected carbon and reliability accounting.

\subsubsection{Evaluation Protocol and Metrics}

All cases are reported under one metric system: first-stage capacities, planning objective, full-year replay objective, runtime, replay curtailment, and replay-feasibility diagnostics. Feasibility is tracked through the replay carbon residual
\begin{align}
r_G(x) = \Big[\mathbb{E}_{\widehat{\mathbb{P}}}[G(x,\xi)] - \bar E^{\mathrm{ann}}/365\Big]_+,
\end{align}
the replay reliability residual
\begin{align}
r_H(x) = \Big[\mathbb{E}_{\widehat{\mathbb{P}}}[H(x,\xi)] - R_{\mathrm{thr}}\Big]_+,
\end{align}
and the number of replay days that violate carbon, reliability, or normal-state no-load-shedding requirements. This common reporting layer is necessary because reduced models may solve quickly while still producing materially different first-stage designs.

\subsubsection{Case 1: Exact Full-Year Benchmark}

The benchmark target is the annual extensive-form problem
\begin{align}
& \text{365 physical days}, \\
& \text{centralized solve}, \\
& \text{second-stage MILP}, \\
& \text{solver: Gurobi}.
\end{align}

This remains the natural reference for all reduced methods. Under the corrected split-cost model, the monolithic 365-day MILP has now been solved to optimality with Gurobi in 2952.65 seconds of wall time. The exact benchmark objective is 8611.54 with zero optimality gap, and the decoded first-stage design is
\begin{align}
P^{\mathrm{pv}}=50,\qquad E^{\mathrm{bess}}=150,\qquad P^{\mathrm{bess}}=50,\qquad P^{\mathrm{grid}}=28,
\end{align}
with carbon-permit purchase 139.93. The exact benchmark also yields zero replay curtailment and annual carbon emission 51076.03~kg, leaving 48923.97~kg slack under the 100000~kg annual cap. All quantitative comparisons below therefore use this exact Case 1 solution rather than the earlier feasible upper-bound replay proxy.

\subsubsection{Case 2: Profile-Based Reduced-Scenario Baseline}

The second test group is the baseline reduced model:
\begin{align}
& \text{$k$-means clustering}, \\
& \text{centralized solve}, \\
& \text{second-stage MILP}, \\
& \text{solver: Gurobi}.
\end{align}

This case quantifies the loss caused by conventional profile-based scenario reduction under the corrected economics. The direct rerun solves quickly at the reduced level, but its first-stage design is highly distorted,
\begin{align}
P^{\mathrm{pv}}=950,\qquad E^{\mathrm{bess}}=400,\qquad P^{\mathrm{bess}}=100,\qquad P^{\mathrm{grid}}=6,
\end{align}
with carbon-permit purchase only 9.52. When that design is replayed on the full-year distribution, the annualized replay objective rises to $7.92\times 10^9$, annual load curtailment reaches 7919.35~kWh, and 55 replay days violate the intended operating requirements. The corrected Case 2 rerun is therefore not merely biased; it is operationally inappropriate as a planning surrogate.

\subsubsection{Case 3: LP Recourse Fidelity Check}

The third test group validates the LP approximation used in the second-stage recourse model. The comparison is performed on a small fixed-plan set consisting of the benchmark plan, the $k$-means baseline plan, and a perturbed benchmark plan. Because the storage state is reset daily, both MILP and LP recourse are evaluated as 365 independent day problems rather than one annual extensive-form model. The resulting error assessment focuses on operating cost, carbon emission, reliability surrogate, and dispatch deviation.

The key error indicators are:
\begin{align}
\varepsilon_Q(x) &:= \left|\mathbb{E}_{\widehat{\mathbb{P}}}[Q^{\mathrm{MILP}}(x,\xi)] - \mathbb{E}_{\widehat{\mathbb{P}}}[Q^{\mathrm{LP}}(x,\xi)]\right|,
\label{eq:10a-lp-cost-error}\\
\varepsilon_G(x) &:= \left|\mathbb{E}_{\widehat{\mathbb{P}}}[G^{\mathrm{MILP}}(x,\xi)] - \mathbb{E}_{\widehat{\mathbb{P}}}[G^{\mathrm{LP}}(x,\xi)]\right|,
\label{eq:10a-lp-carbon-error}\\
\varepsilon_H(x) &:= \left|\mathbb{E}_{\widehat{\mathbb{P}}}[H^{\mathrm{MILP}}(x,\xi)] - \mathbb{E}_{\widehat{\mathbb{P}}}[H^{\mathrm{LP}}(x,\xi)]\right|.
\label{eq:10a-lp-reliability-error}
\end{align}

This experiment validates whether the LP recourse approximation is accurate enough to support the planning-functional reduced-measure construction developed in the theory section above. If \(\varepsilon_Q\), \(\varepsilon_G\), and \(\varepsilon_H\) remain small in the neighborhood of relevant planning solutions, then the LP-based aggregation framework is justified.

For the formal model with SOC-band parameters
\begin{align}
(\theta_1,\theta_2,c^{\mathrm{soc}}_1,c^{\mathrm{soc}}_2)=(0.55,0.70,2.5\times 10^{-5},10^{-4}),
\end{align}
the updated independent-day replay gives the following plan-level errors:
\begin{align}
&\text{benchmark\_full\_milp:} && \varepsilon_Q = 2.04\times 10^{-14}, && \varepsilon_G = 1.14\times 10^{-14}, && \varepsilon_H = 0,\\
&\text{kmeans\_milp:} && \varepsilon_Q = 4.18\times 10^{-9}, && \varepsilon_G = 2.09\times 10^{-15}, && \varepsilon_H = 4.50\times 10^{-15},\\
&\text{benchmark\_grid\_perturbed:} && \varepsilon_Q = 2.07\times 10^{-14}, && \varepsilon_G = 1.25\times 10^{-14}, && \varepsilon_H = 0.
\end{align}
These values remain at numerical-noise level, so the updated Case 3 evidence continues to support LP--MILP equivalence for planning-relevant metrics under the formal SOC-band model. The only visible residual discrepancy is dispatch-level degeneracy: the battery-dispatch \(\ell_1\) distance can still be nonzero in individual scenarios, but these alternative dispatches remain objective-equivalent and carbon-equivalent.

Table~\ref{tab:10a-paper-case3-gaps} and Figure~\ref{fig:10a-case3-lp-gaps} provide the figure-ready summary used in the paper version: both the expected objective gap and the expected carbon / reliability gaps remain at numerical-noise scale across all tested plans.

\input{visualization/microgrid_planning_case10a/tables/table_10a_case3_lp_gaps.tex}

\begin{figure}[t]
\centering
\includegraphics[width=0.78\textwidth]{visualization/microgrid_planning_case10a/figures/fig_10a_case3_lp_gaps.pdf}
\caption{LP--MILP gap comparison for the three fixed-plan validation cases.}
\label{fig:10a-case3-lp-gaps}
\end{figure}

\subsubsection{Case 4: Sequential Scenario-Measure Expansion Results}

The fourth test group evaluates the proposed method under sequential scenario-measure expansion, centralized solution, and a second-stage LP recourse model with grid-trade hull strengthening, battery-throughput regularization, and any enabled formal SOC degradation regularizer. All Case 4 runs reported here use Gurobi as the backend solver.

Because the operational state is reset day by day, Case 4 uses the same independent-day replay semantics as Case 3 when validating each reduced-measure iterate. In the implemented workflow, each iteration first checks whether the planning solution has stabilized relative to the previous iterate; if so, the algorithm stops early and reports planning-solution-stable. Otherwise it replays the current plan on all 365 daily subproblems, inserts the highest replay-violation scenarios into the retained support, then reassigns the entire full-year day set to the nearest current representatives and recomputes representative probabilities from the reassigned masses. Hence the measure update is a global support reallocation aimed at reducing the feature-space Wasserstein discrepancy, not a local probability split on only the newly violated scenarios, and replay-feasible stopping does not trigger a default full-year polish solve. The key verification question is therefore twofold: whether the proposed sequential expansion method can recover a planning solution close to the 365-day benchmark, and whether the monitored support-distance sequence explains the observed stabilization of the planning iterate.

For the calibrated split-cost rerun under this corrected logic, the sequential LP method again terminated after two iterations with replay-feasible stopping. The reduced planning support used 15 scenarios in iteration 1 and 20 scenarios in iteration 2, while the post-update retained support written to disk contained 25 representatives because five additional replay-critical days were inserted after the second replay. The planning and replay objectives moved from
\begin{align}
&5812.17 \text{ and } 3.08\times 10^7 \quad \text{(iteration 1)}
\end{align}
to
\begin{align}
&8576.35 \text{ and } 8615.26 \quad \text{(iteration 2)},
\end{align}
so the replay mismatch again collapsed from a catastrophic first-iteration gap to a small final difference of 38.9. Over the same two iterations, the probability-weighted support distance decreased from 13.2354 to 12.9504, while the maximum covered distance remained 39.2708. Hence the corrected support update still improves the average feature-space approximation even though the worst covered day remains an extreme outlier.

Relative to the earlier annualized-CAPEX smoke rerun under the CAPEX-only regime, the split-cost calibration leaves the physical capacity design unchanged but improves the economic balance. The capex-only smoke run stopped after the same two iterations with
\begin{align}
8576.35 < 8889.07,\qquad 8615.26 < 8930.71,
\end{align}
so the calibrated split-cost accounting lowers both the planning and replay objectives while preserving the same retained support size and replay-feasible stopping pattern. The main numerical shift is in cost allocation: the split-cost run reduces the carbon-permit purchase from 93.47 to 89.35 and moves part of the battery cell cost from first-stage ownership into dispatch-dependent wear charges.

The final support mapping remains structurally reasonable. Its probabilities sum to one up to CSV roundoff and its covered-count total is exactly 365, so the reassignment is mass-conservative. The distribution is concentrated but not collapsed: the ten largest representatives carry about 78.9\% of the annual probability mass, the top fifteen carry about 94.0\%, only three representatives are strict singletons, and eight support points cover at most two days. This is consistent with the intended interpretation of the Wasserstein-oriented update: common day types are merged into a few large representative clusters, while replay-critical edge cases are retained explicitly as low-probability tail support points. In first-stage terms, the final calibrated Case 4 design is
\begin{align}
P^{\mathrm{pv}}=200,\quad P^{\mathrm{bess}}=50,\quad E^{\mathrm{bess}}=150,\quad P^{\mathrm{grid}}=27,
\end{align}
with carbon-permit purchase 89.35.

An additional strict reliability validation was then run with explicit planning-side hard caps
\begin{align}
\bar E^{\mathrm{EENS}}_{\mathrm{day}} = 0, \qquad \bar E^{\mathrm{EENS}}_{\mathrm{ann}} = 0,
\end{align}
implemented through the scenario-level and annual expected-EENS constraints discussed above. Under this hard-cap setting, the sequential Case 4 workflow required three iterations rather than two. The planning and replay objectives moved from
\begin{align}
&5803.42 \text{ and } 3.56\times 10^7 \quad \text{(iteration 1)},\\
&8522.23 \text{ and } 4.59\times 10^7 \quad \text{(iteration 2)},\\
&8523.24 \text{ and } 4.59\times 10^7 \quad \text{(iteration 3)},
\end{align}
while the retained support size expanded from 12 to 14 to 16 and the probability-weighted support distance contracted from 13.6575 to 13.5609 to 13.3577. The final first-stage design remained physically unchanged,
\begin{align}
P^{\mathrm{pv}}=200,\qquad E^{\mathrm{bess}}=150,\qquad P^{\mathrm{bess}}=50,\qquad P^{\mathrm{grid}}=27,
\end{align}
but the associated permit variables became
\begin{align}
\zeta^{\mathrm{CO_2}} = 78.7758,\qquad \zeta^{\mathrm{EENS}} = 0.
\end{align}
Most importantly, the replay-violation report is empty under this strict run: there are no remaining normal-state shedding events, no daily EENS-cap overrun, and no replay carbon-cap violation rows. Hence the planning-side reliability hard caps are not merely diagnostic; they are active model constraints that materially govern the accepted Case 4 iterate.

For compact reporting, Table~\ref{tab:10a-paper-case4-iteration} and Table~\ref{tab:10a-paper-case-summary} now place the corrected Case 4 rerun and the strict hard-cap validation side by side, together with the exact Case 1 benchmark and the corrected Case 2 rerun. This pairing is important: the ordinary corrected rerun shows near-benchmark economic fidelity, while the hard-cap variant shows what happens when replay-feasible planning is additionally required to satisfy explicit planning-side reliability limits rather than only ex post replay diagnostics.

\input{visualization/microgrid_planning_case10a/tables/table_10a_case4_iteration.tex}

\input{visualization/microgrid_planning_case10a/tables/table_10a_case_summary.tex}

\begin{figure}[htbp]
\centering
\includegraphics[width=0.82\textwidth]{visualization/microgrid_planning_case10a/figures/fig_10a_capacity_comparison.pdf}
\caption{Figure-ready first-stage capacity comparison among the exact Case 1 benchmark, the corrected Case 2 rerun, the corrected Case 4 rerun, and the strict Case 4 hard-cap validation.}
\label{fig:10a-capacity-comparison}
\end{figure}

\begin{figure}[htbp]
\centering
\includegraphics[width=0.88\textwidth]{visualization/microgrid_planning_case10a/figures/fig_10a_replay_comparison.pdf}
\caption{Figure-ready comparison of annualized replay objective and replay curtailment. Case 2 remains grossly infeasible under full-year replay, the corrected Case 4 rerun stays economically near-benchmark, and the hard-cap Case 4 validation shows the cost of enforcing zero-EENS planning constraints directly.}
\label{fig:10a-replay-comparison}
\end{figure}

\begin{figure}[htbp]
\centering
\includegraphics[width=0.88\textwidth]{visualization/microgrid_planning_case10a/figures/fig_10a_case4_iteration_progress.pdf}
\caption{Figure-ready Case 4 iteration trajectories for the corrected split-cost rerun and the strict hard-cap validation. The baseline rerun reaches replay-feasible stopping in two iterations, while the hard-cap validation requires a third support expansion step and remains much more expensive in replay objective terms.}
\label{fig:10a-case4-progress}
\end{figure}

The earlier zero-PV conclusion is therefore no longer valid under the corrected economics. The retained representative scenarios still exhibit positive daytime PV capacity factors, so the solar resource was never absent; what changed is the accounting basis. Once PV is charged through annualized CAPEX rather than an effectively unannualized upfront charge, and once battery energy cost is split coherently between fixed ownership and dispatch-dependent wear, a 200 kW PV block becomes economically attractive. Both the corrected rerun and the strict hard-cap validation retain this same physical design, so the qualitative reversal from \(P^{\mathrm{pv}}=0\) to \(P^{\mathrm{pv}}=200\) is driven mainly by the annualization fix, while the reliability hard caps act primarily on the admissible operating policy and permit variables rather than on the installed capacities themselves.

The trust-enabled validation run should be interpreted in this lighter sense. It is not a separate asymptotic solver; it is a two-iteration adaptive support-budget controller layered on top of the same sequential replay workflow. Starting from insertion budget 2, the controller inserts only two replay-critical scenarios after the first iteration, reaches a replay-feasible design already at support size 17 in the second iteration, and then expands the next budget to 3 because the realized merit decrease $1.13\times 10^{10}$ exceeds the predicted decrease $5.01\times 10^8$ by a ratio of 22.48. The final trust-validated design keeps the same physical capacities
\begin{align}
P^{\mathrm{pv}}=200,\qquad E^{\mathrm{bess}}=150,\qquad P^{\mathrm{bess}}=50,\qquad P^{\mathrm{grid}}=27,
\end{align}
but uses lower carbon-permit purchase 86.65 and fewer planning support scenarios than the plain sequential rerun. Because this trust-enabled path also stops after two iterations, its value is chiefly diagnostic: it shows that a smaller insertion budget already captures the replay-critical support corrections needed in this case. Table~\ref{tab:10a-paper-case4-trust} and Figure~\ref{fig:10a-case4-trust} summarize these diagnostics.

\input{visualization/microgrid_planning_case10a/tables/table_10a_case4_trust_validation.tex}

\begin{figure}[htbp]
\centering
\includegraphics[width=0.88\textwidth]{visualization/microgrid_planning_case10a/figures/fig_10a_case4_trust_diagnostics.pdf}
\caption{Two-iteration diagnostic-controller view of the adaptive Case 4 rerun. The merit collapses after the first accepted support update, and the recorded ratio supports expanding the next insertion budget from 2 to 3.}
\label{fig:10a-case4-trust}
\end{figure}

\subsubsection{Experimental Findings}

Under the updated formal-SOC experiments, three conclusions are now supported directly by the corrected result set. First, the old zero-PV narrative is no longer defensible: under annualized CAPEX plus coherent split-cost battery accounting, the validated replay-feasible design is materially solar-positive,
\begin{align}
P^{\mathrm{pv}}=200,\qquad E^{\mathrm{bess}}=150,\qquad P^{\mathrm{bess}}=50,\qquad P^{\mathrm{grid}}=27.
\end{align}
Second, an exact full-year benchmark is now available under the corrected split-cost model: the monolithic Case 1 MILP solves to objective 8611.54 with zero gap, selects the compact design $(50,150,50,28)$ for $(P^{\mathrm{pv}},E^{\mathrm{bess}},P^{\mathrm{bess}},P^{\mathrm{grid}})$, and uses carbon-permit purchase 139.93. Third, the corrected Case 2 rerun confirms that the standard reduced-support baseline is not merely biased but operationally inappropriate: its replay objective explodes to $7.92\times 10^9$, with 7919.35~kWh annual curtailment and 55 violating replay days.

The corrected global support-reassignment logic therefore supports two distinct but complementary conclusions after the economics refactor. On one hand, the ordinary Case 4 rerun reaches replay-feasible stopping in two iterations, collapses the replay objective from $3.08\times 10^7$ to 8615.26, reduces the probability-weighted support distance from 13.2354 to 12.9504, and finishes with only a 3.72 objective gap relative to the exact Case 1 benchmark 8611.54. On the other hand, the new hard-cap validation shows that the same sequential support-expansion machinery can also enforce explicit planning-side reliability limits: with both scenario-level and annual expected-EENS caps fixed at zero, it terminates replay-feasible after three iterations on a 16-point retained support, keeps the same physical design $(200,150,50,27)$, sets the EENS permit to zero, and leaves an empty replay-violation file. The price of that stronger guarantee is economic rather than structural, since the replay objective rises to $4.59\times 10^7$ even though annual replay curtailment remains zero. Case 2, by contrast, still fails by a wide margin under the same replay logic.

The new hard-cap reliability run sharpens this conclusion. With both the scenario-level and annual expected-EENS caps set to zero, the same sequential support-expansion logic still terminates replay-feasible after three iterations on a 16-point retained support, preserves the same physical design $(200,150,50,27)$, and yields an empty replay-violation file. In that stricter setting the replay objective rises to $4.59\times 10^7$, so the hard-cap formulation is no longer near-optimal relative to Case 1 in economic terms; however, it demonstrates that the upgraded planning model can now enforce reliability directly rather than merely reporting it ex post.